
%%%%%%%%%%%%%%%%%%%%%%%%%%%%%%%%%%%%%%%%%%%%%%%%%%%%
%%%%%%%%%%%%%%%%%%%%%%%%%%%%%%%%%%%%%%%%%%%%%%%%%%%%
%%%
%%% For this problem the collaboration and citations statement comes first.
%%%
%%%%%%%%%%%%%%%%%%%%%%%%%%%%%%%%%%%%%%%%%%%%%%%%%%%%
%%%%%%%%%%%%%%%%%%%%%%%%%%%%%%%%%%%%%%%%%%%%%%%%%%%%

\vfill

\paragraph{Problem \arabic{problem} collaboration and citations}
\GradescopeComment{This should be on page 6 (bottom) of your PDF.}

I collaborated on this problem with the following people:
\begin{blue}
  \begin{itemize}
  \item \REPLACEME                    % write "none" if none
  \end{itemize}
\end{blue}

My answers use ideas from the following sources (other than the lecture notes and textbooks): 
\begin{blue}
  \begin{itemize}
  \item \REPLACEME                    % write "none" if none
  \end{itemize}
\end{blue}

\newpage                      % don't delete this line

%%%%%%%%%%%%%%%%%%%%%%%%%%%%%%%%%%%%%%%%%%%%%%%%%%%%
%%%%%%%%%%%%%%%%%%%%%%%%%%%%%%%%%%%%%%%%%%%%%%%%%%%%
%%%
%%%   Problem 4 answer.
%%%
%%%%%%%%%%%%%%%%%%%%%%%%%%%%%%%%%%%%%%%%%%%%%%%%%%%%
%%%%%%%%%%%%%%%%%%%%%%%%%%%%%%%%%%%%%%%%%%%%%%%%%%%%

\paragraph{Problem 4 answer}\GradescopeComment{This should be on page 7 (top) of your PDF.}

%%
%% IF YOU TRY TO PROVE CORRECTNESS, GIVE A PROOF THAT IS AS SIMILAR
%% AS POSSIBLE TO THE ORIGINAL PROOF, BELOW.
%%
%% COLOR ALL TEXT YOU ADD BLUE BY USING THE \BLUE{your text here} MACRO,
%% OR (FOR MULTIPLE LINES) THE \begin{blue} ... \end{blue} ENVIRONMENT
%%

\begin{lemma}\label{lemma: dijkstra invariant}
  In any execution on a graph as described in the problem statement,
  the loop maintains the invariant
  \emph{$d[v] = \dist s  v < \infty$ for all $v\in K$.}
\end{lemma}

\begin{longFormProof}

  \step Consider any execution of the algorithm on some input $(G=(V,E), s)$.

  \step The invariant holds at the start of the first iteration,
  when $K=\{s\}$ and $d[s] = 0 = \dist s s$.
  \\ (Because $G$ has no negative-weight cycles.)

  \begin{block}\label{loop inv}
    
    \step Consider any iteration of the loop such that the invariant holds at the start of the iteration.

    \step Let $K$ and $d$ denote the set $K$ and the array $d$ at the \emph{start} of the iteration.
    
    \step Let $(u',w')$ be the edge chosen in the iteration.
    
    \step To show that the iteration maintains the invariant,
    we'll show $\dist s {w'} = d[u'] + \weight {u', w'}$.

    \smallskip

    \step \underline{First} we show $\dist s {w'} \le d[u'] + \weight {u', w'}$.
    
    \step Since $u'\in K$, by the invariant, $\dist s {u'} = d[u'] < \infty$.

    \step So there exists a path from $s$ to $u'$ of weight $d[u']$.

    \step This path plus the edge $(u', w')$ form a path from $s$ to $w'$ of weight $d[u'] + \weight {u', w'}$.

    \step So there exists a path from $s$ to $w'$ of weight $d[u'] + \weight {u', w'}$.

    \step\label{lt} So $\dist s {w'} \le d[u'] + \weight {u', w'}$, as desired.

    \smallskip
    
    \step \underline{To finish} we'll show $\dist s {w'} \ge d[u'] + \weight {u', w'}$.
    
    \begin{block}\label{P}

      \step Let $P$ be any path from $s$ to $w'$.  We'll show $\weight P \ge d[u'] + \weight {u', w'}$.

      \step $P$ starts in $K$, but ends outside of $K$, so some edge on $P$ leaves $K$.

      \step Let $(x,y)$ be an edge on $P$ that leaves $K$ (so $x\in K$, $y\not\in K$).

      \step Separate the path $P$ around the edge $(x,y)$ into
      the part before, $P_x$, and the part after, $P_y$.

      \step So $P = P_x \cup \{(x,y)\} \cup P_y$ and $P_x$ is a path from $s$ to $x$. Then

      \begin{align*}
        \weight P~
        &=~ \weight {P_x} + \weight {x,y} + \weight {P_y}
        && \comment{(Since $P = P_x \cup \{(x,y)\} \cup P_y$.)}
        \\
        &\ge~ \weight {P_x} + \weight {x,y}
        && \comment{(As all edge weights are non-negative, so $\weight {P_y} \ge 0$.)}
        \\
        &\ge~ \dist s x + \weight {x,y}
        && \comment{(Since $P_x$ is a path from $s$ to $x$.)}
        \\
        &=~ d[x] + \weight {x,y}
        && \comment{(Since $x\in K$ so $d[x] = \dist s x$ by the invariant.)}
        \\
        &\ge~ d[u'] + \weight {u',w'}
        && \comment{(By the alg's choice of $(u', w')$, and $x\in K$, $y\not\in K$)}
      \end{align*}
    \end{block}

    \step By Block~\ref{P}, every path from $s$ to $w'$
    has weight at least $d[u'] + \weight {u',w'}$.

    \step That is, $\dist s {w'} \ge d[u'] + \weight {u',w'}$.

    \smallskip

    \step By this and Step~\ref{lt}, $\dist s {w'} = d[u'] + \weight {u',w'}$.

  \end{block}

  \step The invariant holds initially. By Block~\ref{loop inv}, each iteration maintains it, so it holds throughout.

\end{longFormProof}

\continued

Given that the invariant holds, correctness is easy to verify:  \GradescopeComment{(This should be on Page 8, top.)}

\begin{theorem}
  Given any graph as described in the problem statement,
  the array $d$ returned by \textsf{Dijkstra}'s algorithm
  satisfies $d[v] = \dist s v$ for all $v\in V$.
\end{theorem}
\begin{longFormProof}

  \step Consider the execution of the algorithm on an arbitrary input $(G, s)$.

  \step Consider the time just before Line~4 executes.

  \step By Lemma~\ref{lemma: dijkstra invariant}, at that time, $d[v] = \dist s v$ for each $v\in K$.

  \step By the loop condition, no edges leave $K$, so vertices in $V\setminus K$ are not reachable from $s$.

  \step So each remaining vertex $v\in V\setminus K$ has $\dist s v = \infty$.

  \step So, after Line 4 executes, each vertex $v\in V\setminus K$ also has $d[v] = \dist s v$.

  \step So the distances returned by the algorithm are correct.

\end{longFormProof}


\vfill

%%% NOTE: The collaboration and citation statement for Problem 4 is at the TOP of this file.

\emph{Note: The citation and collaboration statement for Problem 4 is on Page 6 (bottom).}

%%% Local Variables:
%%% mode: latex
%%% TeX-master: "homework_7"
%%% End:
